% Compile with xelatex
\documentclass[11pt,a4paper]{article}

\usepackage{fontspec}
\setmainfont{Malgun Gothic}[Ligatures=TeX]
\setsansfont{Malgun Gothic}[Ligatures=TeX]

\usepackage[top=2.5cm, bottom=2.5cm, left=2.5cm, right=2.5cm]{geometry}
\usepackage{microtype}
\usepackage{amsmath}

\setlength{\parindent}{0pt}
\setlength{\parskip}{6pt}
\pagestyle{empty}

% ─────────────────────────────────────────────────────────────
\begin{document}

\begin{center}
{\large\textbf{Cross-Cultural Perception of Korean Tourist Attractions:\\[3pt]
A Comparative NLP Analysis of Domestic and Foreign Online Reviews}}

\vspace{6pt}

{\normalsize Jorge Adan Soriano Hernandez}\\
{\small Department of Data Science, Hoseo University}

\end{center}

\vspace{6pt}
\noindent\rule{\linewidth}{0.4pt}
\vspace{4pt}

\noindent\textbf{Abstract}

\vspace{2pt}

\noindent
This study investigates how domestic Korean tourists and foreign visitors
perceive the same tourist attractions differently, through a comparative
analysis of online reviews collected from culturally segmented platforms.
A bilingual review corpus is constructed covering major tourist destinations
in Seoul and Busan, including palaces, traditional villages, markets,
entertainment districts, and beaches, selected through stratified sampling
informed by official Korean Tourism Organization (KTO) visitor statistics
and tourism foot-traffic rankings. Each site is represented by balanced
samples of Korean- and English-language reviews. We apply a multi-layered
NLP pipeline comprising: (1)~sentiment analysis using pre-trained classifiers,
with Mann-Whitney U tests for inter-group significance;
(2)~LDA topic modeling to identify dominant themes per language group,
supplemented by TF-IDF keyword contrast and type-token ratio (TTR) analysis;
(3)~cross-lingual semantic similarity via multilingual sentence embeddings
to quantify experiential convergence per site, combined with zero-shot
motivation classification; and (4)~aspect-based sentiment analysis (ABSA)
via instruction-tuned LLMs across five tourism-relevant dimensions
(cultural heritage, accessibility, crowding, food, and visual experience).
Preliminary results suggest that domestic reviewers emphasize practical
concerns such as crowding and revisit logistics, while foreign reviewers
focus on cultural-historical narrative and visual aesthetics. These findings
contribute to the literature on platform-mediated cross-cultural tourism
perception and offer practical implications for destination marketing
organizations seeking to tailor messaging to heterogeneous visitor segments.

\vspace{8pt}
\noindent\rule{\linewidth}{0.4pt}
\vspace{4pt}

\noindent\textbf{초록}

\vspace{2pt}

\noindent
본 연구는 문화적으로 분리된 온라인 리뷰 플랫폼에서 수집한 리뷰를 비교 분석하여,
내국인 관광객과 외국인 방문객이 동일한 한국 관광지를 어떻게 상이하게 인식하는지를
조사한다. 한국관광공사(KTO) 입장객 통계 및 관광 유동인구 순위를 활용한 층화
추출로 서울·부산의 주요 관광지(궁궐, 전통마을, 시장, 유흥가, 해수욕장 등)를
선정하고, 각 장소에 대해 한국어·영어 리뷰를 균형 있게 수집하여 이중언어
코퍼스를 구축한다. 다층적 NLP 파이프라인을 적용하여 (1)~사전 학습된 감성 분류
모델과 Mann-Whitney U 검정을 통한 감성 분석, (2)~LDA 주제 모델링 및 TF-IDF
키워드 대조·유형-토큰 비율(TTR) 분석, (3)~다국어 문장 임베딩을 활용한
교차언어 의미 유사도 측정과 제로샷 방문 동기 분류, (4)~명령 기반 대규모
언어 모델(LLM)을 활용한 다섯 가지 관광 측면(역사·문화, 접근성, 혼잡도, 음식,
시각적 경험)에 대한 측면 기반 감성 분석(ABSA)을 수행한다. 예비 결과에 따르면,
내국인은 혼잡도와 재방문 편의성 등 실용적 측면을, 외국인은 역사·문화적 서사와
시각적 미학을 중시하는 경향이 나타난다. 본 연구는 플랫폼 매개 문화 간 관광
인식 연구에 기여하며, 이질적 방문객 집단에 맞춘 관광 마케팅 전략 수립에
실질적 시사점을 제공한다.

\vspace{6pt}
\noindent\rule{\linewidth}{0.4pt}
\vspace{2pt}

{\small 호서대학교 / Hoseo University \hfill 2026년 10월}

\end{document}
